\section{Data and method}
\label{sec:data}

This section states what the accounts are built from, what they hold fixed, and how a reader
should weigh what comes out of them. It also sets the vocabulary of status labels used in every
table caption that follows.

\subsection{Sources}

Claim levels and holder shares come from the Federal Reserve's financial accounts
\citep{FRBZ1}. The instrument suffix there is not the same across sectors, so matching on the
exact suffix pushes whole holder classes into an unallocated remainder; we match on the
instrument family and take the closest suffix per sector, which cuts that remainder from about a
fifth to under two percent. Every series used is recorded with the publisher's own description in
the replication package.

Household balances come from the Survey of Income and Program Participation \citep{CensusSIPP},
and household composition, earnings and occupation from the American Community Survey
\citep{CensusACS}, with official credit aggregates from the New York Federal Reserve
\citep{NYFedHHDC2026} used to benchmark the survey universe. The wage share of adjusted gross
income, and the realized gains bound that travels with it, are from the Statistics of Income
\citep{IRSSOI}. The distribution of corporate equity is from the distributional financial
accounts \citep{FRBDFA} and behind them the Survey of Consumer Finances \citep{FRBSCF}, and the
housing agency book from the Enterprises' annual filings
\citep{FannieMae10K2025,FreddieMac10K2025}. The side of the accounts that measures claims on AI
capital is built from the filings of nine registrants, with the off-balance-sheet structures
described by \citet{BIS2026} absent from them. The wage bill and federal receipts used as bases
are \result{WageBillBn} and \result{FedReceiptsBn} billion dollars. Units are taken from the
provider at retrieval rather than assumed, and a discontinued series is replaced rather than
carried.

\subsection{What the accounts hold fixed}

Three choices travel with every number in this paper and are worth stating before any of them
appear.

The accounts are a decomposition of a \emph{stock} at a point in time, repeated annually. They
contain no behavioural model, no adoption path and no forecast. Where we report a change between
two dates, it is the difference between two measured decompositions and not an estimate of what
caused it.

Each backing coefficient is a ratio of two measured quantities with a proxy assumption between
them. Table~\ref{tab:bridge} sets out all three for every class, together with an alternative
estimate where one exists, so a reader who rejects a proxy can see exactly which number it
produced.

Every share we report is accompanied by the growth of its own numerator and denominator wherever
the two move differently. A share that looks flat because its denominator grew is a different
fact from a share that is flat because nothing moved, and the long series in
Section~\ref{sec:accounts} contains both.

\subsection{Status labels, and what was independently rebuilt}
\label{sec:rebuild}

Three labels are used throughout and appear in every table caption. \emph{Rebuilt} means the
quantity was reproduced by an independent computational rebuild: a separate run that had no
access to this paper's code and worked from raw data and a published specification alone, inside
a tolerance fixed before the rebuild began. \emph{Measured} means produced by the pipeline,
clearing every applicable bounds check, but not independently rebuilt. \emph{Scenario} means
arithmetic conditional on an assumption stated where the number appears. In prose we give the
label where a reader could otherwise mistake a scenario for a measurement; elsewhere the table
caption carries it.

What a rebuild establishes needs saying plainly, because the word invites more than it should. It
establishes computational reproducibility: that the stated inputs and the stated construction
produce the stated number. It says nothing about whether the construction is the right one, or
whether the economic assumptions behind it hold. A quantity can be rebuilt to six decimal places
and still rest on an assumption we have no evidence for, and several in this paper do.

Three rebuild rounds have been run. Round two scored \result{RepTwoScored} quantities,
\result{RepTwoSealed} of them against expected values fixed in advance; \result{RepTwoFivePct}
came within five percent and \result{RepTwoTolerance} inside the stated tolerance, and
\result{RepTwoMismatch} did not match. Round three covered the two constructions the earlier
rounds had not reached, the debt-only ratio and the assembled capital tax rate, and matched all
\result{RepThreeSealed} quantities, \result{RepThreeRounding} of them within rounding. The blind
was procedural rather than enforced: the rebuild ran on the same machine, with the file of
expected values reachable at a path the specification names, and reported it unopened until the
rebuild was final. The claim is therefore that these quantities were independently rebuilt under
a reported blind protocol, not that they were independently verified.
Appendix~\ref{app:rebuild} gives the round-by-round detail and the causes of the mismatches.

Separately, a standing check states the bound a quantity must respect before it is computed and
tests it afterwards, so that a figure which cannot be right is caught before it is reported
rather than after. Five quantities the accounts can produce are not reported here at all, and
Table~\ref{tab:removed} lists each with the reason.
